% Generated draft: requires mathematical and source review.
\documentclass[11pt]{article}
\usepackage[margin=25mm]{geometry}
\usepackage{fontspec}
\setmainfont{Latin Modern Roman}
\usepackage{amsmath,amssymb,mathtools,unicode-math}
\setmathfont{Latin Modern Math}
\usepackage{enumitem,needspace,xurl,hyperref}
\setlength{\parindent}{0pt}
\setlength{\parskip}{5pt}
\setlength{\emergencystretch}{3em}
\newcommand{\N}{\mathbb{N}}
\newcommand{\Z}{\mathbb{Z}}
\newcommand{\Q}{\mathbb{Q}}
\newcommand{\R}{\mathbb{R}}
\newcommand{\C}{\mathbb{C}}
\newcommand{\F}{\mathbb{F}}
\newcommand{\A}{\mathbb{A}}
\newcommand{\PP}{\mathbb{P}}
\newcommand{\E}{\mathbb{E}}
\newcommand{\eps}{\varepsilon}
\newcommand{\dd}{\,\mathrm{d}}
\begin{document}
\section*{UnsolvedMath ID: 20000747}
\textbf{Source problem code:} catalogue entry
\subsection{Definitions and mathematical statement}
Let \(p_n\) be the \(n\)-th prime. The corrupted source hypothesis is recovered in the supplied research summary as
\[
(\mathrm H)\qquad p_{n+1}-p_n\ll_\varepsilon p_n^\varepsilon
\quad\text{for every fixed }\varepsilon>0.
\]
Precisely, for each \(\varepsilon>0\) there are \(C_\varepsilon<\infty\) and \(n_\varepsilon\) such that the displayed bound holds for every \(n\geq n_\varepsilon\). No constant uniform in \(\varepsilon\), bounded-gap conclusion, or logarithmic gap bound is implicit.

Y. Motohashi asks what analytic information, if any, about \(\zeta(s)\) in \(0<\Re s<1\) can be deduced from \((\mathrm H)\), reversing the usual passage from zeta information to prime-gap upper bounds. The source does not name a particular desired zero-free region or growth estimate, so it is a broad converse question, not an assertion that \((\mathrm H)\) implies the Riemann Hypothesis.

The attached aid stresses that small gaps give one-sided nonemptiness of intervals, not signed prime-counting discrepancy. Its first observation is that the contribution from a fixed zero \(\rho=\beta+i\gamma\), \(\beta<1\), to an explicit-formula interval difference is of size \(O_\rho(yx^{\beta-1})=o(y)\) for \(y\) on a sublinear power scale. This addresses a fixed zero term, not a uniform sum over all zeros or growing heights.

Its second claim is a model obstruction. For every prescribed \(\rho=\beta+i\gamma\) with \(0<\beta<1\) and \(\gamma\ne0\), it reports a bounded sequence \(a_m>0\), \(m\geq2\), such that
\[
D_\rho(s)=\sum_{m\geq2}a_m m^{-s}
\]
has a meromorphic continuation with residues \(+1\) at \(1\) and \(-1\) at \(\rho,\overline\rho\). It reports an associated meromorphic \(Z\) satisfying \(-Z'/Z=D_\rho\), with a pole at \(1\) and simple zeros at the prescribed conjugate pair. This model has dense positive coefficient support but is not the Riemann zeta function and is not supplied with its prime Euler product or other arithmetic structure. Therefore it limits arguments using only support and positivity; it does not disprove a logical implication from actual prime gaps to actual zeta information.

The aid contrasts \((\mathrm H)\) with a genuinely stronger hypothesis involving
\(\psi(x)=\sum_{m\leq x}\Lambda(m)\), where \(\Lambda(p^j)=\log p\) and is zero otherwise. In a nontrivial fixed parameter range \(0<\delta<\eta<1\), assume uniformly for large real \(x\) that
\[
\psi(x+x^\eta)-\psi(x)=x^\eta+O(x^{\eta-\delta}).
\]
It reports that telescoping gives \(\psi(x)=x+O(x^{1-\delta})\), and consequently \(\zeta(s)\) has no zeros in \(\Re s>1-\delta\). This conditional discrepancy implication must not be attributed to the weaker gap hypothesis. The exact broad parameter scope and full proofs of the aid's three claims are not reproduced. They are supplied unreviewed partial obstructions and a stronger-hypothesis converse, not a solution of Motohashi's original question.
\subsection{Short English statement}
Ask what zeta information follows if every prime gap is bounded by \(p_n^\varepsilon\) for each fixed positive \(\varepsilon\). The supplied aid separates interval nonemptiness from signed prime-counting estimates, gives limited model obstructions, and derives a zero-free region only from an explicitly stronger discrepancy assumption.
\subsection{Sources}
\begin{itemize}
\item Y. Motohashi, Problem 11, \emph{Differences between primes and zeta}, December 2005 ARCC workshop \emph{Gaps Between Primes}, notes prepared by Angel Kumchev according to the supplied verification: \url{https://aimath.org/WWN/primegaps/WorkshopProblems.pdf}. The verification reports checking the corrupted hypothesis on the PDF's third page.
\item Local registry \texttt{data/unsolvedmath/problems.json}, permanent ID \(20000747\), exact source code \texttt{catalogue entry}; canonical \texttt{aim-analytic-number-theory-notes.json}, notes[110].
\item The attached research attempt supplies the recovered hypothesis, model construction, and discrepancy converse without a separate complete citation. It is explicitly machine-generated, not peer reviewed, and requires independent verification before citation.
\end{itemize}
\end{document}
